\section{下采样用分辨率换取不变性}
\label{sec:14-Deep-Learning/04-Convolutional-Networks/03}

一段监控视频，摄像头装在杆子上，风一吹画面轻微晃动，每帧的位移不到两个像素。你的分类网络逐帧给出标签，结果标签在相邻帧之间反复跳变：``行人''、``自行车''、``行人''。画面里的人并没有变，变的只有他在传感器上的坐标。

这件事之所以刺眼，是因为你相信卷积网络是平移不变的。它不是。本章第一篇证过的是平移等变——输入移了输出跟着移。要把等变变成不变，网络得在某处主动把位置信息扔掉，而扔的那一下正是跳变的来源。

\subsection{池化与步长卷积在做同一件事}

最大池化取窗口内的最大值，平均池化取均值，步长卷积先把核在每个位置都算一遍再隔 \(s\) 个位置留一个。三者的外形不同，骨架相同：先在局部窗口上算一个响应，再把响应图按 \(s\) 抽稀。

把抽稀单独拎出来记作 \(S_s\)，\((S_s u)_n = u_{sn}\)。步长卷积就是 \(S_s C_w\)，最大池化就是 \(S_s M_k\)（\(M_k\) 是滑动窗口取最大），平均池化就是 \(S_s A_k\)。所有下采样层都是``局部聚合 \(\circ\) 抽稀''这个复合，区别只在聚合算子是线性的还是非线性的。

不变性来自聚合那一步，而且是有条件的。

\begin{proposition}
设 \(M_k\) 为宽度 \(k\) 的滑动最大值。若把输入平移 \(a\) 个位置，且窗口内取到最大值的那个元素在平移后仍落在同一窗口内，则该窗口的输出不变。
\end{proposition}

骨架只有一句：最大值只由``哪个元素最大''和``它在不在窗口里''决定，两个条件都不动，输出就不动。条件对账在于第二个条件。\(\abs{a}\) 相对 \(k\) 越小、最大值元素离窗口边缘越远，条件越容易满足；一旦最大值被平移推出窗口，输出立刻换成窗口里第二大的元素，可以差得很远。所以``池化带来平移不变性''这句话的准确版本是：对小于窗口尺度且不跨越窗口边界的平移近似不变，此外无保证。

代价紧跟着写在维度上。长度 \(T\) 的输入经 \(S_s\) 变成 \(\lceil T/s\rceil\)，映射不是单射，一堆不同的输入被映到同一个输出。这个纤维内部的差异，就是被主动丢掉的信息。想要的那部分叫``小平移不改变结论''，不想要的那部分叫``不再知道东西在哪''——它们是同一个纤维的两种叫法，不可能只留下前者。检测框的定位误差、分割掩码的边界毛刺、小目标在深层特征图上退化成不到一个像素，都是这份账单的分项。

\subsection{抽稀把等变性直接打断}

真正制造帧间跳变的不是聚合，而是抽稀。

\begin{proposition}
记 \((T_a x)_n = x_{n-a}\)。抽稀算子 \(S_s\) 满足 \(S_s T_a = T_{a/s} S_s\) 当且仅当 \(s \mid a\)。
\end{proposition}

证明骨架：\((S_s T_a x)_n = x_{sn-a}\)。要它等于某个 \((T_b S_s x)_n = x_{s(n-b)}\)，需要 \(a = sb\) 对某个整数 \(b\) 成立。条件对账落在整数性上——\(a\) 不是 \(s\) 的倍数时，\(sn-a\) 这个下标集合与 \(sn\) 这个下标集合分属模 \(s\) 的不同剩余类，两次抽稀读的是原信号上互不相交的两组采样点。输出不是原输出的平移，而是一段全新的序列，二者之间没有任何算子层面的关系。

这就是画面晃动一个像素会翻标签的机制。网络里每一个 \(s=2\) 的层，都把平移量按模 2 分成两类；八层下采样累计 stride 32，只有 32 的倍数的平移才被完整保住，其余 31 种相位各自走进不同的采样格。

\begin{center}
\begin{tikzpicture}[x=0.42cm,y=0.30cm,>=stealth]
\begin{scope}[shift={(0,0)}]
  \draw[gray] (-0.8,0) -- (11.8,0);
  \foreach \i in {0,2,4,6,8,10} {\draw[very thick] (\i,0) -- (\i,5); \fill (\i,5) circle (1.5pt);}
  \foreach \i in {1,3,5,7,9,11} {\draw[very thick] (\i,0) -- (\i,1); \fill (\i,1) circle (1.5pt);}
  \foreach \i in {0,2,4,6,8,10} {\draw[dotted] (\i,0) -- (\i,-8.4);}
  \node[font=\scriptsize] at (5.5,7.6) {原信号};
  \draw[->] (5.5,-1.4) -- (5.5,-7.0);
  \node[right,font=\scriptsize,fill=white,inner sep=1.5pt] at (5.7,-4.2) {stride 2};
\end{scope}
\begin{scope}[shift={(0,-8.4)}]
  \draw[gray] (-0.8,0) -- (11.8,0);
  \foreach \i in {0,2,4,6,8,10} {\draw[very thick] (\i,0) -- (\i,5); \fill (\i,5) circle (1.5pt);}
  \node[font=\scriptsize] at (5.5,-2.2) {输出全高};
\end{scope}
\begin{scope}[shift={(15,0)}]
  \draw[gray] (-0.8,0) -- (11.8,0);
  \foreach \i in {0,2,4,6,8,10} {\draw[very thick] (\i,0) -- (\i,1); \fill (\i,1) circle (1.5pt);}
  \foreach \i in {1,3,5,7,9,11} {\draw[very thick] (\i,0) -- (\i,5); \fill (\i,5) circle (1.5pt);}
  \foreach \i in {0,2,4,6,8,10} {\draw[dotted] (\i,0) -- (\i,-8.4);}
  \node[font=\scriptsize] at (5.5,7.6) {右移一个采样点};
  \draw[->] (5.5,-1.4) -- (5.5,-7.0);
  \node[right,font=\scriptsize,fill=white,inner sep=1.5pt] at (5.7,-4.2) {stride 2};
\end{scope}
\begin{scope}[shift={(15,-8.4)}]
  \draw[gray] (-0.8,0) -- (11.8,0);
  \foreach \i in {0,2,4,6,8,10} {\draw[very thick] (\i,0) -- (\i,1); \fill (\i,1) circle (1.5pt);}
  \node[font=\scriptsize] at (5.5,-2.2) {输出全低};
\end{scope}
\end{tikzpicture}
\end{center}

\emph{采样格是钉死在坐标上的，信号一动，同一个格子就从波峰挪到了波谷——输出的差别与平移量毫无比例关系。}

图里的信号在最高频率上振荡，相邻两点一高一低。stride 2 的采样格恰好只碰到波峰，输出是一条高线；把信号右移一个采样点，同一个格子只碰到波谷，输出变成一条低线。两个输出之间不是``平移了一点''，而是完全不同的常数。这正是\hyperref[sec:08-Numerical-Analysis/08-Fourier-and-Spectral-Methods/02]{``采样、混叠与可恢复性''}里讲的那件事：采样率不足以分辨的高频分量，会被折叠成低频，而折叠成什么，取决于采样相位。卷积网络的特征图上到处都是接近 Nyquist 频率的纹理响应——边缘、条纹、网格，恰恰是网络最爱提取的那些。

\subsection{补救的方向都是先低通再抽稀}

采样定理同时给出了处方：抽稀之前把超过新 Nyquist 频率的分量滤掉，混叠就不会发生。工程上的做法是在 stride 层前面插一个固定的低通核（常用二项式模糊核），再做抽稀，即所谓抗混叠下采样。最大池化也可以拆成``窗口取最大（stride 1）\(\to\) 模糊 \(\to\) 抽稀''三步，前两步都是稠密的，只有最后一步降采样。这样改过之后，把输入平移一个像素引起的输出变化会平滑得多，视频上的标签抖动明显减轻。

代价照例存在。低通滤掉的不只是会混叠的伪高频，还包括真实存在的细节——细线、小目标、锐利边界。做分类时这笔交易通常划算，做需要亚像素定位的任务时未必。

另外两条常见路子各有各的边界。数据增强里加随机平移，不改变算子的性质，只是让网络在参数里学出对多种采样相位的容忍，属于用容量买不变性，域外平移仍无保证。全局平均池化把整个空间维压成一个数，确实给出了近似不变，但它把位置信息一次性清零——分类可以，分割和检测不行。

于是就有了下一篇的问题。分割和检测既需要深层那个大感受野带来的语义，又需要浅层没被抽稀掉的空间精度。下采样把分辨率交易出去了，现在得想办法把它换回来——换得回多少，能换回哪一部分，是下一篇要算的账。
